\section{Sílabas abiertas y cerradas}
\textbf{Qué son:} Abiertas terminan en vocal, cerradas en consonante. En neerlandés influyen en la longitud de la vocal.

\textbf{Por qué importan:} Explican ortografía y pronunciación: vocal larga en sílaba abierta, corta en cerrada.

\textbf{Cómo aplicarlo:}
\begin{itemize}[leftmargin=*]
  \item Vocal corta si la sílaba está cerrada: \textit{man, pet, bus}.
  \item Vocal larga si la sílaba está abierta: \textit{ma-ken, le-zen, bo-ten}. Se duplica vocal para mantener sonido largo: \textit{maan}.
  \item Doble consonante para mantener vocal corta antes de sílaba abierta: \textit{man-nen, pet-ten}.
\end{itemize}

\textbf{Ejercicios:}
\begin{itemize}[leftmargin=*]
  \item Separa en sílabas y marca larga/corta: \textit{lopen, stoppen, lezen, bellen}.
  \item Transforma singular a plural y observa cambios de vocal/consonante.
\end{itemize}
